%% =====================================================================
%% NEW CHAPTER: copy this file to Chapters/03Name.tex (number it in reading order) and add
%%     \include{Chapters/03Name}
%% to main.tex, under MAINMATTER. Delete what you don't need below.
%% =====================================================================

%% Chapter title and label (ch:short-name). Page style: mainmatter shows the chapter in the left-hand
%% running head and the section in the right-hand one; use chapteronly for a chapter without sections.
\chapter{Chapter title}\label{ch:name} \pagestyle{mainmatter}

%% Chapter opener (optional): an epigraph and/or a short abstract on the opening page.
%% For no opener, leave both commands with empty braces: \chapterprecishere{} \chapterprecistoc{}
\chapterprecishere{
    % \begin{epigraphs}
    %     \qitem{\textit{Quotation.}}{\headingfont Author}
    % \end{epigraphs}
    % \vspace{0pt plus 1filll} % push the abstract to the bottom of the page
    % \begin{chapterabstract}
    %     What this chapter does and finds, in three to five sentences.
    % \end{chapterabstract}
}
\chapterprecistoc{} % optional one-line summary under the chapter title in the contents

%% Chapter based on a paper (optional): the reference and who did what (CRediT roles, https://credit.niso.org).
% \chapterbasedon{\cite{key}}{A.~Author: conceptualisation, methodology, software, writing -- original draft; B.~Author: supervision, writing -- review \& editing.}

%% Opening paragraph: what the chapter is about and how it is organised.

\section{First section}\label{sec:name-first}

%% Label prefixes: ch: sec: fig: tab: eq: form: alg: def: lem: thm: cor: app:
%% Refer with \cref{...} (\Cref at the start of a sentence).

\subsection{First subsection}\label{sec:name-first-sub}

\section{Conclusions}\label{sec:name-conclusions}
